\chapter{2004年北京卷（春文）}

\section{单选题}

\begin{problem}
在函数 \(y=\sin 2x\)，\(y=\sin x\)，\(y=\cos x\)，\(y=\tan\frac{x}{2}\) 中，最小正周期为 \(\pi\) 的函数是（\quad）

\choices
    {\(y=\sin 2x\)}
    {\(y=\sin x\)}
    {\(y=\cos x\)}
    {\(y=\tan\frac{x}{2}\)}
\end{problem}

\begin{answer}
A
\end{answer}

\begin{solution}
正弦函数 \(y=\sin kx\) 的最小正周期为 \(\dfrac{2\pi}{|k|}\)，所以 \(y=\sin 2x\) 的最小正周期为 \(\pi\)，而 \(y=\sin x\) 和 \(y=\cos x\) 的最小正周期均为 \(2\pi\). 正切函数 \(y=\tan\dfrac{x}{2}\) 的最小正周期为 \(\dfrac{\pi}{1/2}=2\pi\). 因此最小正周期为 \(\pi\) 的函数是 \(y=\sin 2x\)，故选 A.
\end{solution}


\begin{problem}
当 \(m<1\) 时，复数 \(z=2+(m-1)\i\) 在复平面上对应的点位于（\quad）

\choices
    {第一象限}
    {第二象限}
    {第三象限}
    {第四象限}
\end{problem}

\begin{answer}
D
\end{answer}

\begin{solution}
复数 \(z=2+(m-1)\i\) 在复平面上对应的点的横坐标为 \(2>0\)，纵坐标为 \(m-1<0\). 所以该点位于第四象限，故选 D.
\end{solution}


\begin{problem}
双曲线 \(\dfrac{x^{2}}{4}-\dfrac{y^{2}}{9}=1\) 的渐近线方程是（\quad）

\choices
    {\(y=\pm\dfrac{3}{2}x\)}
    {\(y=\pm\dfrac{2}{3}x\)}
    {\(y=\pm\dfrac{9}{4}x\)}
    {\(y=\pm\dfrac{4}{9}x\)}
\end{problem}

\begin{answer}
A
\end{answer}

\begin{solution}
双曲线的标准方程为 \(\dfrac{x^{2}}{a^{2}}-\dfrac{y^{2}}{b^{2}}=1\)，其渐近线方程为 \(y=\pm\dfrac{b}{a}x\). 这里 \(a=2\)，\(b=3\)，所以渐近线方程为 \(y=\pm\dfrac{3}{2}x\)，故选 A.
\end{solution}


\begin{problem}
一个圆锥的侧面积是其底面面积的 \(2\) 倍，则该圆锥的母线与底面所成的角为（\quad）

\choices
    {\(30^{\circ}\)}
    {\(45^{\circ}\)}
    {\(60^{\circ}\)}
    {\(75^{\circ}\)}
\end{problem}

\begin{answer}
C
\end{answer}

\begin{solution}
设圆锥的底面半径为 \(r\)，母线长为 \(l\). 由侧面积是底面面积的 \(2\) 倍，得 \(\pi rl=2\pi r^{2}\)，从而 \(l=2r\). 设母线与底面所成的角为 \(\theta\)，在轴截面的直角三角形中，母线在底面上的投影长为 \(r\)，因此 \(\cos\theta=\dfrac{r}{l}=\dfrac{1}{2}\). 由于 \(0<\theta<90^{\circ}\)，得 \(\theta=60^{\circ}\)，故选 C.
\end{solution}


\begin{problem}
已知 \(\sin(\theta+\pi)<0\)，\(\cos(\theta-\pi)>0\)，则下列不等关系中必定成立的是（\quad）

\choices
    {\(\sin\theta<0\)，\(\cos\theta>0\)}
    {\(\sin\theta>0\)，\(\cos\theta<0\)}
    {\(\sin\theta>0\)，\(\cos\theta>0\)}
    {\(\sin\theta<0\)，\(\cos\theta<0\)}
\end{problem}

\begin{answer}
B
\end{answer}

\begin{solution}
由 \(\sin(\theta+\pi)=-\sin\theta\)，得 \(-\sin\theta<0\)，即 \(\sin\theta>0\). 由 \(\cos(\theta-\pi)=-\cos\theta\)，得 \(-\cos\theta>0\)，即 \(\cos\theta<0\). 所以角 \(\theta\) 的终边位于第二象限，必有 \(\sin\theta>0\) 且 \(\cos\theta<0\)，故选 B.
\end{solution}


\begin{problem}
在抛物线 \(y^{2}=2px\) 上，横坐标为 \(4\) 的点到焦点的距离为 \(5\)，则 \(p\) 的值为（\quad）

\choices
    {\(\dfrac{1}{2}\)}
    {\(1\)}
    {\(2\)}
    {\(4\)}
\end{problem}

\begin{answer}
C
\end{answer}

\begin{solution}
设抛物线上横坐标为 \(4\) 的点为 \(P\left(4,y\right)\). 由 \(y^{2}=2px\)，得 \(y^{2}=8p\). 抛物线 \(y^{2}=2px\) 的焦点为 \(F\left(\dfrac{p}{2},0\right)\)，由 \(PF=5\) 得 \(\left(4-\dfrac{p}{2}\right)^{2}+y^{2}=25\). 代入 \(y^{2}=8p\)，得 \(\left(4-\dfrac{p}{2}\right)^{2}+8p=25\)，即 \(p^{2}+16p-36=0\)，所以 \(p=2\) 或 \(p=-18\). 因为点 \(P\) 在抛物线上且 \(y^{2}=8p\geq0\)，所以 \(p\geq0\)，从而 \(p=2\)，故选 C.
\end{solution}


\begin{problem}
已知 \(a\)，\(b\)，\(c\)，\(d\) 均为实数，有下列命题：
\begin{circlelist}
\item 若 \(ab>0\)，\(bc-ad>0\)，则 \(\dfrac ca-\dfrac db>0\)；
\item 若 \(ab>0\)，\(\dfrac ca-\dfrac db>0\)，则 \(bc-ad>0\)；
\item 若 \(bc-ad>0\)，\(\dfrac ca-\dfrac db>0\)，则 \(ab>0\).
\end{circlelist}
其中正确命题的个数是（\quad）

\choices
    {\(0\)}
    {\(1\)}
    {\(2\)}
    {\(3\)}
\end{problem}

\begin{answer}
D
\end{answer}

\begin{solution}
由于题中出现 \(\dfrac ca\) 和 \(\dfrac db\)，所以 \(a\ne0\)，\(b\ne0\). 记 \(Q=\dfrac ca-\dfrac db\)，则 \(Q=\dfrac{bc-ad}{ab}\).

若 \(ab>0\) 且 \(bc-ad>0\)，则 \(Q>0\)，命题 \(\circled{1}\) 正确. 若 \(ab>0\) 且 \(Q>0\)，则 \(bc-ad=abQ>0\)，命题 \(\circled{2}\) 正确. 若 \(bc-ad>0\) 且 \(Q>0\)，则 \(ab=\dfrac{bc-ad}{Q}>0\)，命题 \(\circled{3}\) 正确. 所以正确命题的个数为 \(3\)，故选 D.
\end{solution}


\begin{problem}
两个完全相同的长方体的长，宽，高分别为 \(5\mathrm{~cm}\)，\(4\mathrm{~cm}\)，\(3\mathrm{~cm}\)，把它们重叠在一起组成一个新长方体，在这些新长方体中，最长的对角线的长度是（\quad）

\choices
    {\(\sqrt{77}\mathrm{~cm}\)}
    {\(7\sqrt2\mathrm{~cm}\)}
    {\(5\sqrt5\mathrm{~cm}\)}
    {\(10\sqrt2\mathrm{~cm}\)}
\end{problem}

\begin{answer}
C
\end{answer}

\begin{solution}
两个相同长方体沿某一个方向拼接后仍为长方体，因而新长方体的三条棱长可能为 \(10\mathrm{~cm}\)，\(4\mathrm{~cm}\)，\(3\mathrm{~cm}\)，或 \(5\mathrm{~cm}\)，\(8\mathrm{~cm}\)，\(3\mathrm{~cm}\)，或 \(5\mathrm{~cm}\)，\(4\mathrm{~cm}\)，\(6\mathrm{~cm}\). 相应的对角线长的平方分别为 \(10^{2}+4^{2}+3^{2}=125\)，\(5^{2}+8^{2}+3^{2}=98\)，\(5^{2}+4^{2}+6^{2}=77\). 最大值为 \(125\)，所以最长对角线长为 \(\sqrt{125}=5\sqrt5\mathrm{~cm}\)，故选 C.
\end{solution}


\begin{problem}
在 \(100\) 件产品中有 \(6\) 件次品，现从中任取 \(3\) 件产品，至少有 \(1\) 件次品的不同取法的种数是（\quad）

\choices
    {\(\mathrm C_{6}^{1}\mathrm C_{94}^{2}\)}
    {\(\mathrm C_{6}^{1}\mathrm C_{99}^{2}\)}
    {\(\mathrm A_{100}^{3}-\mathrm A_{94}^{3}\)}
    {\(\mathrm C_{100}^{3}-\mathrm C_{94}^{3}\)}
\end{problem}

\begin{answer}
D
\end{answer}

\begin{solution}
从 \(100\) 件产品中任取 \(3\) 件的取法总数为 \(\mathrm C_{100}^{3}\). 没有次品时，必须从 \(94\) 件正品中取出 \(3\) 件，取法数为 \(\mathrm C_{94}^{3}\). 因此至少取到 \(1\) 件次品的取法数为 \(\mathrm C_{100}^{3}-\mathrm C_{94}^{3}\)，故选 D.
\end{solution}


\begin{problem}
期中考试以后，班长算出了全班 \(40\) 个人数学成绩的平均分为 \(M\)，如果把 \(M\) 当成一个同学的分数，与原来的 \(40\) 个分数一起，算出这 \(41\) 个分数的平均值为 \(N\)，那么 \(M:N\) 为（\quad）

\choices
    {\(\dfrac{40}{41}\)}
    {\(1\)}
    {\(\dfrac{41}{40}\)}
    {\(2\)}
\end{problem}

\begin{answer}
B
\end{answer}

\begin{solution}
设原来 \(40\) 个分数的总和为 \(S\)，则 \(S=40M\). 加入一个分数 \(M\) 后，\(41\) 个分数的总和为 \(S+M=41M\)，所以 \(N=\dfrac{41M}{41}=M\). 因而 \(M:N=1\)，故选 B.
\end{solution}


\section{填空题}

\begin{problem}
直线 \(x-\sqrt3y+a=0\)（\(a\) 为常实数）的倾斜角的大小是\fillinblank{}.
\end{problem}

\begin{answer}
\(30^{\circ}\)
\end{answer}

\begin{solution}
将直线方程化为斜截式，得 \(y=\dfrac{1}{\sqrt3}x+\dfrac{a}{\sqrt3}\)，所以直线的斜率为 \(\dfrac{1}{\sqrt3}\). 设倾斜角为 \(\alpha\)，则 \(\tan\alpha=\dfrac{1}{\sqrt3}\)，且 \(0\leq\alpha<180^{\circ}\)，故 \(\alpha=30^{\circ}\).
\end{solution}

\begin{problem}
\(\dfrac{\sin(\alpha+30^{\circ})-\sin(\alpha-30^{\circ})}{\cos\alpha}\) 的值为\fillinblank{}.
\end{problem}

\begin{answer}
\(1\)
\end{answer}

\begin{solution}
由正弦的差化积公式，得 \(\sin(\alpha+30^{\circ})-\sin(\alpha-30^{\circ})=2\cos\alpha\sin30^{\circ}=\cos\alpha\). 原式有意义时 \(\cos\alpha\ne0\)，所以原式等于 \(1\).
\end{solution}

\begin{problem}
若 \(f^{-1}(x)\) 为函数 \(f(x)=\lg(x-1)\) 的反函数，则 \(f^{-1}(x)\) 的值域是\fillinblank{}.
\end{problem}

\begin{answer}
\((1,+\infty)\)
\end{answer}

\begin{solution}
设 \(y=\lg(x-1)\)，则函数的定义域为 \(x>1\). 由对数定义得 \(10^{y}=x-1\)，所以 \(x=10^{y}+1\)，从而反函数为 \(f^{-1}(x)=10^{x}+1\). 因为 \(10^{x}>0\)，且当 \(x\) 取遍实数时 \(10^{x}\) 取遍正数，所以 \(f^{-1}(x)\) 的值域为 \((1,+\infty)\).
\end{solution}

\begin{problem}
若直线 \(mx+ny-3=0\) 与圆 \(x^{2}+y^{2}=3\) 没有公共点，则 \(m\)，\(n\) 满足的关系式为\fillinblank{}；以 \((m,n)\) 为点 \(P\) 的坐标，过点 \(P\) 的一条直线与椭圆 \(\dfrac{x^{2}}7+\dfrac{y^{2}}3=1\) 的公共点有\fillinblank{}个.
\end{problem}

\begin{answer}
\(m^{2}+n^{2}<3\)；\(2\)
\end{answer}

\begin{solution}
由于题目已说明 \(mx+ny-3=0\) 是一条直线，所以 \(m\)，\(n\) 不同时为 \(0\). 圆 \(x^{2}+y^{2}=3\) 的圆心为原点，半径为 \(\sqrt3\). 直线 \(mx+ny-3=0\) 到原点的距离为 \(\dfrac{3}{\sqrt{m^{2}+n^{2}}}\)，因此直线与圆没有公共点的充要条件是
\(\dfrac{3}{\sqrt{m^{2}+n^{2}}}>\sqrt3\)，即 \(m^{2}+n^{2}<3\).

结合 \(m^{2}+n^{2}>0\)，可知实际上 \(0<m^{2}+n^{2}<3\)，填空中的关系式通常记为 \(m^{2}+n^{2}<3\).

由此，点 \(P\left(m,n\right)\) 满足 \(\dfrac{m^{2}}{7}+\dfrac{n^{2}}{3}\leq\dfrac{m^{2}+n^{2}}{3}<1\)，所以 \(P\) 在椭圆 \(\dfrac{x^{2}}7+\dfrac{y^{2}}3=1\) 的内部. 任取过 \(P\) 的直线，设其参数方程为 \(\left(x,y\right)=\left(m,n\right)+t\left(u,v\right)\)，其中 \(\left(u,v\right)\ne\left(0,0\right)\). 代入椭圆方程，得到关于 \(t\) 的二次方程 \(At^{2}+2Bt+C=0\)，其中 \(A=\dfrac{u^{2}}7+\dfrac{v^{2}}3>0\)，\(C=\dfrac{m^{2}}7+\dfrac{n^{2}}3-1<0\). 因此两根的乘积 \(\dfrac{C}{A}<0\)，方程有两个异号的实根，过点 \(P\) 的直线与椭圆有两个公共点.
\end{solution}

\section{解答题}

\begin{problem}
解不等式：\(\sqrt{2x-1}>x-2\).
\end{problem}

\begin{solution}
因为根式有意义，必须有 \(2x-1\geq0\)，即 \(x\geq\dfrac{1}{2}\).

当 \(\dfrac{1}{2}\leq x<2\) 时，\(\left(x-2\right)<0\)，而 \(\sqrt{2x-1}\geq0\)，所以不等式恒成立.

当 \(x\geq2\) 时，不等式两边均为非负数，可以平方，得 \(2x-1>\left(x-2\right)^{2}\)，即 \(x^{2}-6x+5<0\)，所以 \(\left(x-1\right)\left(x-5\right)<0\)，从而 \(1<x<5\). 结合 \(x\geq2\)，得 \(2\leq x<5\).

综上，原不等式的解集为 \(\left[\dfrac{1}{2},5\right)\).
\end{solution}

\begin{problem}
在 \(\triangle ABC\) 中，\(a\)，\(b\)，\(c\) 分别是 \(\angle A\)，\(\angle B\)，\(\angle C\) 的对边长，已知 \(a\)，\(b\)，\(c\) 成等比数列，且 \(a^{2}-c^{2}=ac-bc\)，求 \(\angle A\) 的大小及 \(\dfrac{b\sin B}{c}\) 的值.
\end{problem}

\begin{solution}
因为 \(a\)，\(b\)，\(c\) 成等比数列且均为三角形的边长，所以 \(b^{2}=ac\). 设 \(t=\sqrt{\dfrac{a}{c}}>0\)，则 \(a=ct^{2}\)，\(b=ct\). 将其代入已知等式，得 \(c^{2}(t^{4}-1)=c^{2}(t^{2}-t)\)，即 \(t^{4}-t^{2}+t-1=0\).

因 \(t^{4}-t^{2}+t-1=(t-1)(t^{3}+t^{2}+1)\)，且 \(t^{3}+t^{2}+1>0\)，所以 \(t=1\)，从而 \(a=b=c\). 因此 \(\triangle ABC\) 为等边三角形，\(\angle A=60^{\circ}\)，且 \(B=60^{\circ}\). 于是 \(\dfrac{b\sin B}{c}=\sin60^{\circ}=\dfrac{\sqrt3}{2}\).
\end{solution}

\begin{problem}
如图，四棱锥 \(S-ABCD\) 的底面是边长为 \(1\) 的正方形，\(SD\) 垂直于底面 \(ABCD\)，\(SB=\sqrt3\).

\begin{enumerate}
\item 求证：\(BC\perp SC\)；
\item 求面 \(ASD\) 与面 \(BSC\) 所成二面角的大小.
\end{enumerate}

\texfigure[max width=0.42\linewidth]{img_repaint/2004/beijing_liberal_spring/q17_fig1.tex}
\FloatBarrier
\end{problem}

\begin{solution}
建立空间直角坐标系，取 \(D\) 为原点，令 \(DA\)，\(DC\)，\(DS\) 分别为 \(x\)，\(y\)，\(z\) 轴正方向. 由题意可设 \(A\left(1,0,0\right)\)，\(C\left(0,1,0\right)\)，\(B\left(1,1,0\right)\)，\(S\left(0,0,h\right)\). 由 \(SB=\sqrt3\)，得 \(SB^{2}=1^{2}+1^{2}+h^{2}=3\)，所以 \(h=1\).

于是 \(\overrightarrow{BC}=\left(-1,0,0\right)\)，\(\overrightarrow{SC}=\left(0,1,-1\right)\)，从而 \(\overrightarrow{BC}\cdot\overrightarrow{SC}=0\)，故 \(BC\perp SC\).

平面 \(ASD\) 的方程为 \(y=0\)，可取法向量 \(\bs{n}_{1}=\left(0,1,0\right)\). 平面 \(BSC\) 过 \(B\left(1,1,0\right)\)，\(S\left(0,0,1\right)\)，\(C\left(0,1,0\right)\)，其方程为 \(y+z=1\)，可取法向量 \(\bs{n}_{2}=\left(0,1,1\right)\). 设两平面的所成二面角为 \(\varphi\)，则 \(\cos\varphi=\dfrac{|\bs{n}_{1}\cdot\bs{n}_{2}|}{|\bs{n}_{1}||\bs{n}_{2}|}=\dfrac{1}{\sqrt2}\)，所以 \(\varphi=45^{\circ}\).
\end{solution}

\begin{problem}
\(2003\) 年 \(10\) 月 \(15\) 日 \(9\) 时，“神舟”五号载人飞船发射升空，于 \(9\) 时 \(9\) 分 \(50\) 秒准确进入预定轨道，开始巡天飞行. 该轨道是以地球的中心 \(F_{2}\) 为一个焦点的椭圆. 选取坐标系如图所示，椭圆中心在原点. 近地点 \(A\) 距地面 \(200\mathrm{~km}\)，远地点 \(B\) 距地面 \(350\mathrm{~km}\). 已知地球半径 \(R=6371\mathrm{~km}\).

\begin{enumerate}
\item 求飞船飞行的椭圆轨道的方程；
\item 飞船绕地球飞行了 \(14\) 圈后，于 \(16\) 日 \(5\) 时 \(59\) 分返回舱与推进舱分离，结束巡天飞行，飞船共巡天飞行了约 \(6\times10^{5}\mathrm{~km}\)，问飞船巡天飞行的平均速度是多少 \(\mathrm{km/s}\)？（结果精确到 \(1\mathrm{~km/s}\)）
\end{enumerate}

\texfigure[max width=0.46\linewidth]{img_repaint/2004/beijing_liberal_spring/q18_fig1.tex}
\FloatBarrier
\end{problem}

\begin{solution}
\begin{enumerate}
\item 设椭圆的长半轴为 \(a\)，焦半距为 \(c\)，短半轴为 \(b\). 由图可知，\(F_{2}\) 是右焦点，\(A\) 是近地点，\(B\) 是远地点，所以 \(AF_{2}=a-c=6371+200=6571\)，\(BF_{2}=a+c=6371+350=6721\). 因此 \(2a=6571+6721=13292\)，\(2c=6721-6571=150\)，从而 \(a=6646\)，\(c=75\).

又 \(b^{2}=a^{2}-c^{2}=6646^{2}-75^{2}=44163691\)，故椭圆轨道方程为
\[
\dfrac{x^{2}}{6646^{2}}+\dfrac{y^{2}}{44163691}=1
\]

\item 从15日9时9分50秒到16日5时59分，共飞行 \(20\) 小时 \(49\) 分 \(10\) 秒，即 \(20\times3600+49\times60+10=74950\mathrm{~s}\). 因此平均速度为 \(v=\dfrac{6\times10^{5}}{74950}\mathrm{~km/s}\approx8.005\mathrm{~km/s}\)，按题意精确到 \(1\mathrm{~km/s}\) 得平均速度约为 \(8\mathrm{~km/s}\).
\end{enumerate}
\end{solution}

\begin{problem}
某服装厂生产一种服装，每件服装的成本为 \(40\) 元，出厂单价定为 \(60\) 元，该厂为鼓励销售商订购，决定当一次订购量超过 \(100\) 件时，每多订购一件，订购的全部服装的出厂单价就降低 \(0.02\) 元. 根据市场调查，经销商一次订购量不会超过 \(500\) 件.

\begin{enumerate}
\item 设一次订购量为 \(x\) 件，服装的实际出厂单价为 \(P\) 元，写出函数 \(P=f(x)\) 的表达式；
\item 当销售商一次订购 \(400\) 件服装时，该服装厂获得的利润是多少元？（服装厂售出一件服装的\(\text{利润} = \text{实际出厂单价} - \text{成本}\)）
\end{enumerate}
\end{problem}

\begin{solution}
\begin{enumerate}
\item 当 \(1\leq x\leq100\) 时，订购量没有超过 \(100\) 件，实际出厂单价仍为 \(60\) 元. 当 \(100<x\leq500\) 时，超过 \(100\) 件的部分有 \(x-100\) 件，订购的全部服装每件降低 \(0.02\left(x-100\right)\) 元，所以实际出厂单价为 \(60-0.02\left(x-100\right)=62-0.02x\) 元. 因此
\[
P=f(x)=
\begin{cases}
60, & 1\leq x\leq100,\\
62-0.02x, & 100<x\leq500.
\end{cases}
\]

\item 订购 \(400\) 件时，超过 \(100\) 件的部分为 \(400-100=300\) 件，所以每件服装的实际出厂单价为 \(60-0.02\times300=54\) 元. 每件服装的利润为 \(54-40=14\) 元，故服装厂获得的总利润为 \(400\times14=5600\) 元.
\end{enumerate}
\end{solution}

\begin{problem}
下表给出一个“等差数阵”.

\begin{center}
\renewcommand{\arraystretch}{1.15}
\begin{tabular}{|c|c|c|c|c|c|c|c|}
\hline
\(4\)&\(7\)&\(\left(\ \right)\)&\(\left(\ \right)\)&\(\left(\ \right)\)&\(\cdots\)&\(a_{1j}\)&\(\cdots\)\\ \hline
\(7\)&\(12\)&\(\left(\ \right)\)&\(\left(\ \right)\)&\(\left(\ \right)\)&\(\cdots\)&\(a_{2j}\)&\(\cdots\)\\ \hline
\(\left(\ \right)\)&\(\left(\ \right)\)&\(\left(\ \right)\)&\(\left(\ \right)\)&\(\left(\ \right)\)&\(\cdots\)&\(a_{3j}\)&\(\cdots\)\\ \hline
\(\left(\ \right)\)&\(\left(\ \right)\)&\(\left(\ \right)\)&\(\left(\ \right)\)&\(\left(\ \right)\)&\(\cdots\)&\(a_{4j}\)&\(\cdots\)\\ \hline
\(\vdots\)&\(\vdots\)&\(\vdots\)&\(\vdots\)&\(\vdots\)&\(\vdots\)&\(\vdots\)&\(\vdots\)\\ \hline
\(a_{i1}\)&\(a_{i2}\)&\(a_{i3}\)&\(a_{i4}\)&\(a_{i5}\)&\(\cdots\)&\(a_{ij}\)&\(\cdots\)\\ \hline
\(\vdots\)&\(\vdots\)&\(\vdots\)&\(\vdots\)&\(\vdots\)&\(\vdots\)&\(\vdots\)&\(\vdots\)\\ \hline
\end{tabular}
\end{center}
其中每行，每列都是等差数列，\(a_{ij}\) 表示位于第 \(i\) 行第 \(j\) 列的数.

\begin{enumerate}
\item 写出 \(a_{45}\) 的值；
\item 写出 \(a_{ij}\) 的计算公式以及 \(2008\) 这个数在等差数阵中所在的位置.
\end{enumerate}
\end{problem}

\begin{solution}
\begin{enumerate}
\item 记第 \(i\) 行的公差为 \(h_i\)，第 \(j\) 列的公差为 \(v_j\). 由表中前两行、前两列可得 \(h_1=7-4=3\)，\(h_2=12-7=5\)，\(v_1=7-4=3\)，\(v_2=12-7=5\). 对任意 \(i\)，\(j\)，从 \(a_{ij}\) 先向右再向下，或先向下再向右，都可到达 \(a_{i+1,j+1}\)，所以 \(h_i+v_{j+1}=v_j+h_{i+1}\)，即 \(h_{i+1}-h_i=v_{j+1}-v_j\). 因此这些相邻公差的增量都是 \(2\)，从而 \(h_i=3+2\left(i-1\right)=2i+1\).

又 \(a_{i1}=4+3\left(i-1\right)=3i+1\)，所以 \(a_{ij}=a_{i1}+\left(j-1\right)h_i=3i+1+\left(j-1\right)\left(2i+1\right)=2ij+i+j\). 于是 \(a_{45}=2\times4\times5+4+5=49\).

\item 由上式，令 \(a_{ij}=2008\)，得
\[
2ij+i+j=2008\Longleftrightarrow(2i+1)(2j+1)=4017=3\times13\times103
\]
由于 \(i\)，\(j\) 为正整数，\(2i+1\)，\(2j+1\) 是大于 \(1\) 的奇数. 因而所有可能的有序因数对为
\[
\begin{split}
(2i+1,2j+1)\in\{&(3,1339),(13,309),(39,103),\\
&(103,39),(309,13),(1339,3)\}.
\end{split}
\]
所以
\[
(i,j)\in\{(1,669),(6,154),(19,51),(51,19),(154,6),(669,1)\}
\]
即 \(2008\) 位于第 \(1\) 行第 \(669\) 列、第 \(6\) 行第 \(154\) 列、第 \(19\) 行第 \(51\) 列、第 \(51\) 行第 \(19\) 列、第 \(154\) 行第 \(6\) 列和第 \(669\) 行第 \(1\) 列.
\end{enumerate}
\end{solution}
